\begin{named}{Proposition~\ref{prop:research}(d): cheap research}
Suppose the classes satisfy the class ranking, the prior is \(1/2\), \(0<\rho<1\), \(0<\beta<1/2\),
\(a,k,c_H>0\) and \(0<\lb<1\). An indifferent investor may mix over its best orders and an indifferent high-cost
challenger prepares. Full preparation at the purchase also gives this outcome when an indifferent high-cost
challenger may prepare with any probability.
Let \((F_r)_{r\in[r_-,r_+]}\) be a path of continuous incumbent laws, strengthening the incumbent in first-order
stochastic dominance with every other primitive fixed, with \(\DT>0\), \(g_H>g_L\), continuous
\(\DT,\tau\), and \(\tau>1/2\). Suppose the path reaches the cap spread, and let
\(r_{\lb}\) be the first strength at which it does, with the spread strictly below that value at \(r_-\).
On \([r_-,r_{\lb}]\), suppose \(\tau<1\). Write
\[
 J(r)=\DT(r)[1-\tau(r)],\qquad
 \varphi_a^*(r)=J(r)-k+(\rho J(r)-k)^+-2a\lambda_E(\tau(r)),
\]
\[
 \lambda_E(\tau)=\frac{\beta(2\tau-1)}{\beta+(1-2\beta)(1-\tau)}.
\]
If \(\varphi_a^*\) is strictly increasing on \([r_-,r_{\lb}]\) and
\(\varphi_a^*(r_-)<0<\varphi_a^*(r_{\lb})\), recruiting research is an equilibrium
exactly on
\[
 [r_R,\ \max\{r\in[r_-,r_+]:\tau(r)\le M_{\lb}\}],
\]
where \(r_R\in(r_-,r_{\lb})\) is the unique zero of \(\varphi_a^*\).
Research is below its cap on \([r_R,r_{\lb})\) and at the cap thereafter on this interval.
On a path with continuously differentiable \(\DT,\tau\) on \([r_-,r_{\lb}]\), the cheap-research condition
\begin{equation}\tag{CR}\label{eq:a5-cheap}
 [1+\rho\1\{\rho J(r)>k\}]J'(r)
 -2a\frac{\beta\tau'(r)}{[\beta+(1-2\beta)(1-\tau(r))]^2}>0
 \quad\text{for almost every }r\in[r_-,r_{\lb}]
\end{equation}
makes \(\varphi_a^*\) strictly increasing there.
\end{named}

\begin{proof}
For a fixed \(\DT>2k\), the function
\[
 \DT m_\lambda-k+(\rho\DT m_\lambda-k)^+-2a\lambda
\]
is continuous and strictly decreasing in \(\lambda\), positive at zero and negative at one. Its unique zero is
\(\lambda_R(\DT)\). At \(\lambda=\lambda_E(\tau)\), the function equals \(\varphi_a^*\), because
\(m_{\lambda_E}=1-\tau\). Therefore, when \(1/2<\tau<1\), the purchase posterior reaches the threshold at the
capped root exactly when
\[
 \varphi_a^*\ge0\quad\hbox{and}\quad\lambda_E(\tau)\le\lb.
\]
A nonnegative margin implies \(J-k+(\rho J-k)^+\ge2a\lambda_E>0\). If \(J\le k\), its left side is nonpositive,
so \(J>k\), which gives \(\DT>2k\) since \(1-\tau<1/2\). Thus these two inequalities also imply that the root
exists. Conversely, the root comparison proves their necessity.

Continuity, strict increase and the endpoint signs of \(\varphi_a^*\) yield exactly one zero \(r_R\), with a
negative margin below it and a positive margin above it.

At \(r_{\lb}\), continuity makes the spread equal to the cap spread, so \(\lambda_R=\lb\). The positive margin there gives
\(\lambda_E(\tau(r_{\lb}))<\lb\), or \(\tau(r_{\lb})<M_{\lb}\). Since \(\tau\) is nondecreasing, the cap
posterior reaches the threshold throughout \([r_-,r_{\lb}]\). On that interval, the exact root comparison
therefore makes recruitment equivalent to \(\varphi_a^*\ge0\), or \(r\ge r_R\). Below \(r_{\lb}\), the spread
is strictly below the cap spread, so research is below its cap. From \(r_{\lb}\) on, research optimality at the
cap persists because \(\DT\) is nondecreasing. The cap result gives exactly the upper interval in the statement,
including its last point and every threshold plateau within it.

When \(\DT,\tau\) are continuously differentiable, \(\varphi_a^*\) is absolutely continuous on
\([r_-,r_{\lb}]\). Away from \(\rho J=k\) its derivative is the expression in \eqref{eq:a5-cheap}.
On the level set \(\rho J=k\), \(J'=0\) almost everywhere, so that expression is its derivative almost
everywhere there as well. Integrating the strictly positive derivative over any nondegenerate subinterval
proves that \(\varphi_a^*\) rises strictly.
\end{proof}

\paragraph{The research margin.}
Where \(\tau'>0\), condition~\eqref{eq:a5-cheap} says that
\[
 a<\frac{[1+\rho\1\{\rho J>k\}]J'}{2\lambda_E'(\tau)\tau'}.
\]
Where \(\tau'=0\), it says \(J'>0\). Before a sure sale pays, the derivative of the margin is
\(J'-2a\lambda_E'(\tau)\tau'\); once it pays, the derivative is
\((1+\rho)J'-2a\lambda_E'(\tau)\tau'\). The margin is continuous at \(\rho J=k\); at a crossing into
\(\rho J>k\), its derivative jumps up by \(\rho J'\). The condition applies where the recruiting interval
starts with a strictly profitable sure sale.

\paragraph{Cheap research in the example.}
On \([r_0,r_N]\), a sale at the prior earns at most its trading cost, and \(\tau>1/2\), so \(\rho J<k\) and
\(\varphi_a^*=\varphi_a=J-k-2a\lambda_E\). In the example \(r_{\lb}<r_N<\reg{res_bound_strength}\).
For \(r\in[r_0,\reg{res_bound_strength}]\), differentiating the auction formulas gives
\[
 \DT'(r)\ge\reg{res_bound_dDelta},\quad
 1-\tau(r)\ge\reg{res_bound_tau74},\quad
 \DT(r)\le\reg{res_bound_Delta74},\quad
 \tau'(r)\le\reg{res_bound_taup}.
\]
Here \(\DT'\), \(\DT\) and \(\tau\) rise, while \(-g_L'\) and \(\gamma=g_H-g_L\) fall. The formula
\(\tau'=(-g_L'+\tau\DT')/\gamma\) gives the upper bound
\[
 \tau'(r)\le
 \frac{-g_L'(r_0)+\tau(\reg{res_bound_strength})\DT'(\reg{res_bound_strength})}
 {\gamma(\reg{res_bound_strength})}
 <\reg{res_bound_taup}.
\]
Also \(\lambda_E'(\tau)\) rises with \(\tau\). Substitution in
\(\varphi_a'=\DT'(1-\tau)-[\DT+2a\lambda_E'(\tau)]\tau'\) therefore gives
\[
 J'(r)\ge \reg{res_bound_dDelta}\cdot\reg{res_bound_tau74}
 -\reg{res_bound_Delta74}\cdot\reg{res_bound_taup}>0,\qquad
 \varphi_a'(r)\ge\reg{res_phip_lb}>0
 \quad\text{at }a=\reg{res_a}.
\]
The same bound is positive for \(a<a_{\rm mono}\approx\reg{res_amono}\).
The endpoint values are
\[
 \varphi_a^*(r_0)=\reg{res_phi_r0_frac}<0,
 \qquad \varphi_a^*(r_{\lb})\approx\reg{res_phi_rcap}>0.
\]
Hence recruiting research begins at the unique zero \(r_R\approx\reg{res_rR}\), reaches its cap at
\(r_{\lb}\approx\reg{res_rcap}\), and continues to \(r_C\approx\reg{res_rC}\).
By Lemma~\ref{lem:a4-path}(b), \(J\) rises on the whole path. Since \(\lambda_E\) also rises, for
\(r\in[\reg{res_bound_strength},r_2]\),
\[
 \varphi_a(r)\ge J(\reg{res_bound_strength})-k-2a\lambda_E(r_2)>0,
 \qquad J(\reg{res_bound_strength})=\reg{res_J74_frac}.
\]
Thus \(r_R\) is the unique zero of \(\varphi_a\) on \([r_0,r_2]\).
